% FINEX crayon notebook theme; the upstream Eisvogel template stays unmodified.
\usepackage{tikz}
\usetikzlibrary{calc,shapes.geometric,decorations.pathmorphing,shadows}
\usepackage{etoolbox}
\usepackage{fvextra}
\usepackage{newunicodechar}
\usepackage{needspace}
% Eisvogel loads mdframed after this header. TikZ gives the notes real rounded
% corners and a quiet paper shadow while retaining mdframed's page breaks.
\PassOptionsToPackage{framemethod=tikz}{mdframed}
\definecolor{FinexInk}{HTML}{3F3A36}
\definecolor{FinexMuted}{HTML}{917D6D}
\definecolor{FinexCream}{HTML}{FFF8EF}
\definecolor{FinexPeach}{HTML}{F7D9C4}
\definecolor{FinexClay}{HTML}{B85F48}
\definecolor{FinexSage}{HTML}{C9D9BC}
\definecolor{FinexSageDark}{HTML}{53785C}
\definecolor{FinexLilac}{HTML}{DCD2E9}
\definecolor{FinexHoney}{HTML}{F1CE78}
\definecolor{FinexBlue}{HTML}{5578C4}
\definecolor{FinexMint}{HTML}{A9DCC5}
\definecolor{FinexPaper}{HTML}{FFFCF6}
\definecolor{FinexPlum}{HTML}{806297}
\colorlet{FinexAccent}{FinexSageDark}
\colorlet{FinexWash}{FinexMint!24!white}

% Local handwritten faces for headings. Portable fonts remain available as a
% fallback through the same TeX installation/bundle used by the body fonts.
\IfFontExistsTF{KaiTi}{%
  \newCJKfontfamily\FinexHandCJK[AutoFakeBold=1.5]{KaiTi}%
}{\newCJKfontfamily\FinexHandCJK[AutoFakeBold=1.5]{FandolKai-Regular.otf}}
\IfFontExistsTF{Comic Sans MS}{%
  \newfontfamily\FinexHandLatin{Comic Sans MS}%
}{\newfontfamily\FinexHandLatin{texgyreheros-italic.otf}}
\newcommand{\FinexHand}{\FinexHandLatin\FinexHandCJK}
\DeclareRobustCommand{\FinexHandText}[1]{{\FinexHand #1}}
% A rounded face makes the large cover title feel playful rather than calligraphic.
\IfFontExistsTF{YouYuan}{%
  \newCJKfontfamily\FinexCoverCJK[AutoFakeBold=2]{YouYuan}%
}{\newCJKfontfamily\FinexCoverCJK[AutoFakeBold=2]{FandolHei-Regular.otf}}
\DeclareRobustCommand{\FinexCoverText}[1]{{\FinexHandLatin\FinexCoverCJK #1}}
\AtBeginEnvironment{titlepage}{%
  \let\textsf\FinexCoverText
  % Keep the stronger cover hierarchy local to the title page.
  \renewcommand{\huge}{\fontsize{36}{44}\selectfont}% company name
  \renewcommand{\Large}{\fontsize{20}{28}\selectfont}% subtitle, author and date
}
\tikzset{finex pencil/.style={decorate,decoration={random steps,segment length=3pt,amplitude=.35pt},line cap=round,line join=round}}

% A restrained four-color rhythm keeps the chapter tabs from feeling repetitive.
\newcommand{\FinexSelectAccent}{%
  \pgfmathtruncatemacro{\FinexPaletteIndex}{mod(max(\value{section}-1,0),4)}%
  \ifcase\FinexPaletteIndex
    \colorlet{FinexAccent}{FinexSageDark}\colorlet{FinexWash}{FinexMint!24!white}%
  \or
    \colorlet{FinexAccent}{FinexBlue}\colorlet{FinexWash}{FinexBlue!10!white}%
  \or
    \colorlet{FinexAccent}{FinexPlum}\colorlet{FinexWash}{FinexLilac!43!white}%
  \or
    \colorlet{FinexAccent}{FinexClay}\colorlet{FinexWash}{FinexPeach!44!white}%
  \fi}
\newcommand{\FinexSpark}[1]{%
  \draw[#1,line width=.65pt,line cap=round]
    (-.09,0)--(.09,0) (0,-.12)--(0,.12);
  \draw[#1,line width=.45pt,line cap=round]
    (-.05,-.065)--(.05,.065) (-.05,.065)--(.05,-.065);}
\newcommand{\FinexListBullet}{\tikz[baseline=-.45ex,x=1em,y=1em]{%
  \fill[FinexBlue!25] (0,0)--(.26,.13)--(.05,.30)--(-.16,.13)--cycle;
  \draw[FinexBlue,line width=.45pt,line join=round]
    (0,0)--(.26,.13)--(.05,.30)--(-.16,.13)--cycle;}}
\newcommand{\FinexNoteTitle}[2]{%
  \tikz[baseline=(note.base)]{%
    \node[anchor=base west,fill=#1!12!white,rounded corners=2pt,
      text=#1,inner xsep=8pt,inner ysep=3pt,font=\FinexHand\bfseries]
      (note) {#2};
    \draw[#1!55,line width=.65pt,line cap=round]
      ($(note.north east)+(-7pt,2pt)$)--($(note.north east)+(-3pt,-2pt)$);
    \draw[#1!55,line width=.65pt,line cap=round]
      ($(note.north east)+(-3pt,3pt)$)--($(note.north east)+(1pt,-1pt)$);}}

% KOMA retains the real heading, numbering, bookmarks and multi-line wrapping.
% The decoration is local to the title, so longer reports need no fixed layout.
\let\FinexOriginalSectionLinesFormat\sectionlinesformat
\renewcommand*{\sectionlinesformat}[4]{%
  \ifstrequal{#1}{section}{%
    \begingroup
    \FinexSelectAccent
    \begin{tikzpicture}[baseline=(heading.base)]
      \node[anchor=base west,align=left,fill=FinexWash,rounded corners=5pt,
        text width=\dimexpr\linewidth-30pt\relax,
        inner xsep=8pt,inner ysep=5pt] (heading)
        {#3\textcolor{FinexAccent}{#4}};
      \draw[FinexAccent!48,line width=.65pt,finex pencil]
        ($(heading.south west)+(12pt,-1.5pt)$)--
        ($(heading.south west)+(65pt,-1pt)$);
      \begin{scope}[shift={($(heading.east)+(-2pt,1pt)$)},x=1cm,y=1cm]
        \FinexSpark{FinexAccent!70}
      \end{scope}
    \end{tikzpicture}%
    \endgroup
  }{\FinexOriginalSectionLinesFormat{#1}{#2}{#3}{#4}}}

% Common report markers use portable vector/math glyphs instead of emoji fonts.
\newcommand{\FinexWarning}{\tikz[baseline=-.2ex,x=1em,y=1em]{%
  \draw[FinexClay,fill=FinexHoney!35,rounded corners=.5pt,line width=.5pt]
    (0,0)--(.85,0)--(.425,.74)--cycle;
  \node[text=FinexInk,font=\tiny\bfseries,inner sep=0pt] at (.425,.27) {!};}}
\newcommand{\FinexFace}[1]{\tikz[baseline=-.4ex,x=1em,y=1em]{%
  \fill[FinexHoney] (0,0) circle (.42);
  \fill[FinexInk] (-.14,.08) circle (.025);
  \fill[FinexInk] (.14,.08) circle (.025);
  \draw[FinexInk,line width=.4pt] (-.15,-.14) .. controls (0,#1) .. (.15,-.14);}}
\newcommand{\FinexStar}[1]{\tikz[baseline=-.3ex]{%
  \node[star,star points=5,star point ratio=2.2,draw=FinexClay!75,
    fill=#1,minimum size=1em,inner sep=0pt,line width=.4pt] {};}}
\newunicodechar{★}{\FinexStar{FinexHoney}}
\newunicodechar{☆}{\FinexStar{white}}
\newunicodechar{✅}{\textcolor{FinexSageDark}{\ensuremath{\checkmark}}}
\newunicodechar{✓}{\textcolor{FinexSageDark}{\ensuremath{\checkmark}}}
\newunicodechar{❌}{\textcolor{FinexClay}{\ensuremath{\times}}}
\newunicodechar{✗}{\textcolor{FinexClay}{\ensuremath{\times}}}
\newunicodechar{⚠}{\FinexWarning}
\newunicodechar{🟢}{\textcolor{FinexSageDark}{\ensuremath{\bullet}}}
\newunicodechar{🔴}{\textcolor{FinexClay}{\ensuremath{\bullet}}}
\newunicodechar{🟡}{\textcolor{FinexHoney!65!FinexInk}{\ensuremath{\bullet}}}
\newunicodechar{📈}{\textcolor{FinexSageDark}{\ensuremath{\nearrow}}}
\newunicodechar{😟}{\FinexFace{-.02}}
\newunicodechar{😐}{\FinexFace{-.14}}
\newunicodechar{😊}{\FinexFace{-.3}}
\newunicodechar{🤝}{\textcolor{FinexSageDark}{[协作]}}
\newunicodechar{🔥}{\textcolor{FinexClay}{[热度]}}
% Ignore the emoji/text presentation selector (U+FE0F), not the marker itself.
\newunicodechar{️}{}

\AtBeginDocument{%
  \color{FinexInk}
  \definecolor{heading-color}{HTML}{B85F48}
  \definecolor{blockquote-border}{HTML}{A9DCC5}
  \definecolor{blockquote-text}{HTML}{705B4B}
  \definecolor{table-row-color}{HTML}{FFF3E5}
  \definecolor{table-rule-color}{HTML}{89A4D9}
  \addtokomafont{section}{\FinexHand\color{FinexClay}\Large}
  \addtokomafont{subsection}{\FinexHand\color{FinexBlue}\normalsize}
  \addtokomafont{subsubsection}{\FinexHand\color{FinexSageDark}\normalsize}
  \setkomafont{pageheadfoot}{\FinexHand\scriptsize\color{FinexMuted}}
  \KOMAoptions{headsepline=false,footsepline=false}
  \pagecolor{FinexPaper}
  % Eisvogel loads mdframed after this header, so define panels at document start.
  \newmdenv[backgroundcolor=FinexMint!12!white,linecolor=FinexSageDark!33,
    linewidth=.65pt,roundcorner=7pt,skipabove=10pt,skipbelow=12pt,
    shadow=true,shadowsize=2pt,shadowcolor=FinexInk!9!white,
    innertopmargin=9pt,innerbottommargin=11pt,innerleftmargin=13pt,innerrightmargin=13pt,
    frametitle={\FinexNoteTitle{FinexSageDark}{研究速览}},frametitleaboveskip=9pt,
    frametitlebelowskip=7pt]{FinexSummary}
  \newmdenv[backgroundcolor=FinexPeach!21!white,linecolor=FinexClay!38,
    linewidth=.65pt,roundcorner=7pt,skipabove=10pt,skipbelow=12pt,
    shadow=true,shadowsize=2pt,shadowcolor=FinexInk!9!white,
    innertopmargin=9pt,innerbottommargin=11pt,innerleftmargin=13pt,innerrightmargin=13pt,
    frametitle={\FinexNoteTitle{FinexClay}{风险备忘}},frametitleaboveskip=9pt,
    frametitlebelowskip=7pt]{FinexRisk}
  \AtBeginEnvironment{longtable}{%
    \rowcolors{2}{FinexHoney!13!white}{FinexMint!11!white}%
    \renewcommand{\arraystretch}{1.4}%
    \arrayrulecolor{FinexBlue!55!white}}
  \renewcommand{\labelitemi}{\FinexListBullet}
  \renewcommand{\labelitemii}{\textcolor{FinexSageDark}{\ensuremath{\circ}}}
  \renewcommand{\labelenumi}{\textcolor{FinexBlue}{\FinexHandLatin\theenumi.}}
  \renewcommand{\tightlist}{\setlength{\itemsep}{3pt}\setlength{\parskip}{0pt}}
  \renewcommand*{\pagemark}{%
    \tikz[baseline=(page.base)]{%
      \node[fill=FinexHoney!33!white,rounded corners=2pt,rotate=-4,
        inner xsep=5pt,inner ysep=2pt,font=\FinexHandLatin\scriptsize]
        (page) {\thepage};}}
  \ofoot*{\pagemark}
  \renewcommand*{\sectionformat}{%
    \tikz[baseline=(badge.base)]{%
      \node[fill=FinexAccent,text=white,rounded corners=2pt,rotate=-4,
        inner xsep=7pt,inner ysep=3pt,font=\FinexHandLatin\small\bfseries]
        (badge){\thesection};
      \draw[FinexAccent!60,line width=.6pt,line cap=round]
        ($(badge.south west)+(1pt,-2pt)$)--($(badge.south east)+(1pt,-1pt)$);
    }\enskip}
  \RedeclareSectionCommand[beforeskip=1.35\baselineskip,afterskip=.65\baselineskip]{section}
  \RedeclareSectionCommand[beforeskip=1.0\baselineskip,afterskip=.3\baselineskip]{subsection}
  % Keep a chapter tab with enough of its first paragraph/list to start reading.
  \pretocmd{\section}{\Needspace{7\baselineskip}}{}{%
    \PackageWarning{finex-warm}{Could not reserve chapter opening space}}
  \setlength{\emergencystretch}{3em}
  \fvset{breaklines=true,breakanywhere=true,fontsize=\small}
}

% Hand-drawn notebook details sit in the margins, outside the reading area.
% Count physical pages because Eisvogel resets its page counter after the cover.
\newcounter{FinexPhysicalPage}
\AddToHook{shipout/before}{\stepcounter{FinexPhysicalPage}}
\AddToHook{shipout/background}{%
  \ifnum\value{FinexPhysicalPage}>1
    \begin{tikzpicture}[remember picture,overlay,x=1cm,y=1cm]
      \coordinate (origin) at (current page.south west);
      \begin{scope}[shift={(origin)}]
        \pgfmathsetseed{\value{FinexPhysicalPage}}
        \draw[FinexMint!50!white,line width=2.1pt,finex pencil]
          (2.3,27.83)--(18.65,27.80);
        \draw[FinexBlue!65!white,line width=.55pt,finex pencil]
          (2.3,27.82)--(18.68,27.81);
        \draw[FinexPeach!70!white,line width=1.9pt,finex pencil]
          (2.3,1.22)--(18.7,1.24);
        \draw[FinexClay!50!white,line width=.4pt,finex pencil]
          (2.3,1.21)--(18.7,1.23);
        % Small crayon hatching at the outer edges, away from tables and prose.
        \foreach \offset in {0,.11,.22,.33}{
          \draw[FinexLilac!60!white,line width=2.4pt,finex pencil]
            (.45+\offset,25.95)--(.83+\offset,26.35);
          \draw[FinexMint!60!white,line width=2.2pt,finex pencil]
            (19.7+\offset,3.0)--(20.02+\offset,3.52);
        }
        % Alternate a tiny sun, a star trail and a leaf sprig across pages.
        \pgfmathtruncatemacro{\FinexDoodleIndex}{mod(\value{FinexPhysicalPage},3)}
        \ifcase\FinexDoodleIndex
          % Keep tiny closed curves undecorated: PGF random steps can overflow.
          \draw[FinexHoney!85!FinexClay,fill=FinexHoney!20!white,
            line width=.7pt] (19.9,26.8) ellipse (.16 and .14);
          \foreach \angle in {0,60,...,300}{%
            \draw[FinexHoney!85!FinexClay,line width=.7pt,line cap=round]
              (19.9,26.8)++(\angle:.23)--++(\angle:.10);}
        \or
          \node[star,star points=5,star point ratio=2.2,draw=FinexBlue!65,
            fill=FinexBlue!10!white,line width=.6pt,rotate=13,
            minimum size=.38cm,inner sep=0pt] at (19.85,26.85) {};
          \draw[FinexBlue!45,line width=.6pt,dash pattern=on 1pt off 2pt]
            (19.85,26.55) .. controls (20.25,26.25) and (19.55,26.15) .. (19.9,25.85);
        \or
          \draw[FinexSageDark!65,line width=.7pt,line cap=round]
            (19.7,26.45) .. controls (20.0,26.7) .. (20.05,27.05);
          \filldraw[draw=FinexSageDark!55,fill=FinexMint!45!white,line width=.5pt]
            (19.86,26.63) .. controls (19.48,26.68) and (19.55,26.93) .. (19.92,26.82)
            .. controls (19.99,26.54) and (20.24,26.56) .. (20.01,26.91);
        \fi
        \begin{scope}[shift={(1.05,2.25)},rotate=-12,x=1cm,y=1cm]
          \draw[FinexPlum!55,line width=.7pt,finex pencil]
            (-.18,-.15)--(.16,-.12)--(.19,.17)--(-.14,.20)--cycle;
          \draw[FinexPlum!45,line width=.5pt,line cap=round]
            (-.08,.08)--(.10,.08) (-.08,0)--(.07,0);
        \end{scope}
      \end{scope}
    \end{tikzpicture}%
  \fi
}

% The cover artwork is painted by Eisvogel. These labels stay sharp and editable.
\AddToHook{shipout/foreground}{%
  \ifnum\value{FinexPhysicalPage}=1
    \begin{tikzpicture}[remember picture,overlay]
      \node[anchor=west,text=FinexInk,font=\FinexHandLatin\bfseries\fontsize{32}{38}\selectfont]
        at ($(current page.south west)+(3.5cm,25.9cm)$) {FINEX};
      \node[anchor=west,text=FinexMuted,font=\FinexHandLatin\fontsize{13}{17}\selectfont]
        at ($(current page.south west)+(3.5cm,25.1cm)$) {RESEARCH, WITH A LITTLE WARMTH};
      \draw[FinexBlue!70,line width=1.8pt,finex pencil]
        ($(current page.south west)+(3.5cm,24.65cm)$)--($(current page.south west)+(7.7cm,24.7cm)$);
    \end{tikzpicture}%
  \fi
}
